\documentclass[final,t]{beamer}
\mode<presentation>
\usetheme{default}
\usepackage[orientation=portrait,size=a0,scale=1.4]{beamerposter}
\usepackage{graphicx}
\usepackage{booktabs}
\usepackage{tikz}
\usepackage{amsmath,amsthm,amssymb}

% Colors
\definecolor{umnmaroon}{RGB}{122,0,25}
\definecolor{umngold}{RGB}{255,204,51}
\setbeamercolor{block title}{fg=white,bg=umnmaroon}
\setbeamercolor{block body}{fg=black,bg=white}
\setbeamercolor{block alerted title}{fg=white,bg=umngold}
\setbeamercolor{block alerted body}{fg=black,bg=white}

% Layout
\setbeamertemplate{navigation symbols}{}
\setbeamertemplate{caption}[numbered]

\title{\huge Compiling Structured Operational Semantics to Executable OCaml Code}
\author{Linglong Meng}
\institute{
\large
    University of Minnesota, Twin Cities}
\date{}

\begin{document}

\begin{frame}[t]
\begin{columns}[t]

% Header spanning both columns
\begin{column}{\paperwidth}
\vskip1cm
\begin{beamercolorbox}[wd=\paperwidth,center]{title}
\vskip1cm
\usebeamerfont{title}\inserttitle\par
\vskip1cm
\usebeamerfont{author}\insertauthor\par
\vskip0.5cm
\usebeamerfont{institute}\insertinstitute\par
\vskip1cm
\end{beamercolorbox}

\begin{beamercolorbox}[wd=\paperwidth,center,sep=1cm]{block body}
\large
\textbf{2025 Midwest Programming Languages Summit}\\
December 13, 2025 | University of Minnesota, Twin Cities
\end{beamercolorbox}
\vskip1cm
\end{column}

\end{columns}

\begin{columns}[t]

% Left Column
\begin{column}{.48\paperwidth}

\begin{block}{Introduction}
\large
When we write down specifications for a programming langauges 
and reason about it, we usually write down on fancy paper specs, 
then in someway turn into implementation. 
However, here we are trying to build a compiler 
that can directly translate into OCaml code. 

\end{block}

\vskip2ex

\begin{block}{Background}
\large
The compiler is build on a existed compiler 
build upon by UMN graudted Ph.D Dawn Michaelson. 
The Compiler Sterling can generate Prolog (as 
one of its many other backends like extensibella 
but that's for reasoning framework). 
\end{block}

\vskip2ex

\begin{block}{Spec Language Overview}
\large
\begin{minipage}[t]{0.48\textwidth}
\textbf{Constructor declaration:}
\begin{align*}
\tau ::= &\; \alpha \\
       | &\; \beta \\
       | &\; \kappa(\gamma) \\
       | &\; \delta(\epsilon, \zeta, \eta)
\end{align*}
\end{minipage}
\hfill
\begin{minipage}[t]{0.48\textwidth}
\textbf{Overall structure:}
\[
\frac{\begin{array}{c}
\langle\text{premise 1}\rangle \\
\vdots \\
\langle\text{premise n}\rangle
\end{array}}{\langle\text{conclusion}\rangle} \; [\text{rule name}]
\]
\end{minipage}
Judgement declaration, you can see it as function signature. 


\end{block}

\vskip2ex

\begin{block}{Example}
\large
Provide a concrete example of your work here.
\end{block}

\end{column}

% Right Column
\begin{column}{.48\paperwidth}

\begin{block}{Results}
\large
Present your main results or findings here.
\end{block}

\vskip2ex

\begin{block}{Implementation}
\large
Discuss implementation details, such as:
\begin{itemize}
\item OCaml code generation strategy
\item Key design decisions
\item Performance considerations
\end{itemize}
\end{block}

\vskip2ex

\begin{block}{Evaluation}
\large
Include evaluation metrics, benchmarks, or case studies.
\end{block}

\vskip2ex

\begin{block}{Conclusion \& Future Work}
\large
Summarize your contributions and outline future directions.
\end{block}

\vskip2ex

\begin{block}{References}
\normalsize
\begin{thebibliography}{99}
\bibitem{ref1} Author1. Title1. Conference/Journal, Year.
\bibitem{ref2} Author2. Title2. Conference/Journal, Year.
\end{thebibliography}
\end{block}

\end{column}

\end{columns}

\vskip1cm

% Footer
\begin{columns}[t]
\begin{column}{\paperwidth}
\begin{beamercolorbox}[wd=\paperwidth,center,sep=0.5cm]{block body}
\small
\textit{hello}\\
Contact: meng0181@umn.edu
\end{beamercolorbox}
\end{column}
\end{columns}

\end{frame}

\end{document}